\chapter{Исследовательская часть}

\section{Оценка временной эффективности}

Мерой временной эффективности рассматриваемых алгоритмов выступает количество операций сравнения искомого значения $x$ с элементами массива $A$.

\subsection{Поиск полным перебором}

Лучший случай реализуется, если искомый элемент расположен на первой позиции массива. В этом случае алгоритм выполняет единственное сравнение, что даёт:
\begin{equation}
	T_{\min}(N) = 1 = O(1).
\end{equation}

Худший случай реализуется, если искомое значение отсутствует в массиве либо находится на последней позиции. Тогда требуется последовательно проверить все $N$ элементов:
\begin{equation}
	T_{\max}(N) = N = O(N).
\end{equation}

\subsection{Бинарный поиск}

Принцип бинарного поиска основан на последовательном уменьшении области поиска вдвое. На каждом шаге алгоритм сравнивает искомое значение со средним элементом текущего обрабатываемой части.

Обозначим через $N_i$ размер обрабатываемой части после $i$-й итерации, при этом $N_0 = N$. Рекуррентное соотношение имеет вид:
\begin{align}
	N_1 &= \frac{N}{2}, \notag \\
	N_2 &= \frac{N}{2^2}, \notag \\
	&\;\;\vdots \notag \\
	N_i &= \frac{N}{2^i}.
\end{align}

Процесс завершается при достижении обрабатываемой части размера единица: $N / 2^i = 1$, откуда получаем $i = \log_2 N$. Таким образом, количество итераций ограничено логарифмической функцией, а итоговое число сравнений определяется структурой конкретного варианта реализации.

\subsubsection{Нерекурсивный вариант с отложенным выходом}

В данной реализации цикл продолжается до тех пор, пока границы диапазона не совпадут ($l = r$). На каждой итерации выполняется единственное сравнение $A[m] < x$, определяющее направление уменьшения обрабатываемой части. После завершения цикла производится финальная проверка $A[l] == x$.

Поскольку досрочное завершение отсутствует, лучший и худший случаи совпадают. Общее число сравнений составляет:
\begin{equation}
	T_{\min}(N) = T_{\max}(N) = \lfloor\log_2 N\rfloor + 1 = O(\log_2 N).
\end{equation}

\subsubsection{Нерекурсивный вариант с досрочным выходом}

На каждой итерации данного алгоритма выполняются два сравнения: сначала проверяется равенство $A[m] == x$, затем --- порядок $A[m] < x$. При обнаружении совпадения поиск немедленно завершается.

Лучший случай --- искомый элемент совпадает со средним элементом исходного массива. В этом случае достаточно одного сравнения:
\begin{equation}
	T_{\min}(N) = 1 = O(1).
\end{equation}

Худший случай --- элемент отсутствует или расположен на границе обрабатываемой части. Тогда на каждой из $\lceil\log_2(N+1)\rceil$ итераций выполняются оба сравнения:
\begin{equation}
	T_{\max}(N) = 2\lceil\log_2(N+1)\rceil = O(\log_2 N).
\end{equation}

\subsubsection{Рекурсивный вариант}

Рекурсивная реализация эквивалентна нерекурсивному варианту с досрочным выходом по числу выполняемых сравнений. Каждый рекурсивный вызов содержит две операции сравнения ($A[m] == x$ и $A[m] < x$), а глубина рекурсии ограничена величиной $\lceil\log_2(N+1)\rceil$:
\begin{align}
	T_{\min}(N) &= 1 = O(1), \notag \\
	T_{\max}(N) &= 2\lceil\log_2(N+1)\rceil = O(\log_2 N).
\end{align}

\section{Оценка затрачиваемой памяти}

Оценка памяти производится только для дополнительных структур данных, не включая входной массив.

\subsection{Линейный поиск}

Алгоритм оперирует двумя служебными переменными:
\begin{itemize}
	\item $i$ --- счётчик текущего элемента;
	\item $index$ --- переменная для хранения результата.
\end{itemize}

Функция вызывается однократно, дополнительные затраты памяти не зависят от размера входных данных: $O(1)$.

\subsection{Нерекурсивный бинарный поиск}

Оба нерекурсивных варианта используют три вспомогательные переменные:
\begin{itemize}
	\item $l$ --- левая граница обрабатываемой части;
	\item $r$ --- правая граница обрабатываемой части;
	\item $m$ --- индекс среднего элемента обрабатываемой части.
\end{itemize}

Поскольку количество переменных фиксировано и не растёт с увеличением $N$, дополнительная память составляет $O(1)$.

\subsection{Рекурсивный бинарный поиск}

Каждый рекурсивный вызов создаёт в стеке новый кадр, содержащий параметры функции (ссылку на массив $A$, значение $x$, границы $l$ и $r$), локальную переменную $m$, адрес возврата и служебные данные. Обозначим объём одного кадра через $M_s$.

Максимальная глубина рекурсии определяется скоростью сужения диапазона и равна $\lceil\log_2(N+1)\rceil$. Суммарные затраты памяти в наихудшем случае:
\begin{equation}
	M_{\text{rec}}(N) = \lceil\log_2(N+1)\rceil \cdot M_s = O(\log_2 N).
\end{equation}

Следовательно, рекурсивная реализация требует логарифмического объёма дополнительной памяти, в то время как нерекурсивные варианты обходятся константным.

\section{Сравнительный анализ}

Результаты теоретического сравнения алгоритмов представлены в таблице~\ref{tbl:comparison}, где ЛС и ХС обозначают наилучший и наихудший случаи соответственно.

\begin{table}[H]
\caption{Сравнение алгоритмов поиска}
\label{tbl:comparison}
\centering
\small
\begin{tblr}{
  width = \linewidth,
  colspec = {
    X[2.2,c,m]
    X[0.9,c,m]
    X[0.9,c,m]
    X[2.2,c,m]
    X[2.4,c,m]
  },
  hlines, vlines,
  row{1,2} = {font=\bfseries},
}
\SetCell[r=2]{c} {\diagbox[width=3.6cm, height=2.2cm]{Алгоритм}{Критерии}}
  & \SetCell[c=2]{c} {Оценка временной\\эффективности как\\функции от числа\\необходимых\\сравнений}
  & 
  & \SetCell[r=2]{c} {Требования,\\налагаемые\\алгоритмом}
  & \SetCell[r=2]{c} {Оценка памяти,\\затрачиваемой\\алгоритмом сверх\\входных данных} \\
  & ЛС & ХС & & \\
Поиск полным перебором
  & $1$
  & $N$
  & Нет
  & {$i \to O(1)$,\\$1$ вызов функции} \\
{Нерекурсивный бинарный\\поиск (отложенный выход)}
  & $1 + \log_2 N$
  & $1 + \log_2 N$
  & \SetCell[r=3]{c} {Массив\\предварительно\\отсортирован\\и поддерживается\\отсортированным}
  & \SetCell[r=2]{c} {$l, r, m \to O(1)$,\\$1$ вызов функции} \\
{Нерекурсивный бинарный\\поиск (предв. выход)}
  & $1$
  & $2 \cdot (1 + \log_2 N)$ 
  & 
  & \\
Рекурсивный бинарный поиск
  & $1$
  & $2 \cdot (1 + \log_2 N)$ 
  & 
  & {$\log_2 N \cdot M_s \to O(\log_2 N)$,\\$\log_2 N$ вызовов} \\
\end{tblr}
\end{table}

\subsection*{Выводы по таблице~\ref{tbl:comparison}:}
\begin{itemize}
	\item Линейный поиск в наихудшем случае требует $N$ сравнений, однако не предъявляет требований к упорядоченности данных. Его применение оправдано для малых массивов или неотсортированных коллекций.
	
	\item Бинарный поиск обеспечивает логарифмическую сложность $O(\log_2 N)$, но требует предварительную сортировку массива и поддержание его в отсортированном состоянии.
	
	\item Вариант реализации с досрочным выходом достигает константной сложности $O(1)$ в лучшем случае, однако в худшем случае выполняет вдвое больше сравнений по сравнению с вариантом с отложенным выходом.
	
	\item Рекурсивная реализация демонстрирует ту же временную эффективность, что и нерекурсивный вариант с досрочным выходом, но затрачивает $O(\log_2 N)$ дополнительной памяти на организацию стека вызовов. Следует отметить дополнительное ограничение рекурсивной реализации: глубина рекурсии ограничена размером стека вызовов, предоставляемого операционной системой. При больших значениях $N$ существует теоретическая опасность переполнения стека, чего лишены нерекурсивные варианты.
	
	\item Для поиска в неотсортированном массиве оптимален линейный поиск. Для отсортированных данных при $N \gg 1$ предпочтителен нерекурсивный бинарный поиск с отложенным выходом: он гарантирует минимальное число сравнений в наихудшем случае ($\approx \log_2 N$), обладает предсказуемым временем выполнения и не расходует память на рекурсивные вызовы.
\end{itemize}

\section{Вывод}

В результате исследовательской части были проведены оценки временной эффективности и затрачиваемой памяти для каждого алгоритма поиска. Был проведён сравнительный анализ полученных результатов.